% !TeX root = ../常微分方程笔记.tex
% 由主文档通过 \input 载入；请编译主文档。
\setcounter{chapter}{13} % 保留原书章号。
\chapter{数学建模与综合专题}
\label{chap:ode-modeling}

一个完整模型由状态变量、演化规律、参数、初始数据与适用范围组成。数学分析应回到原问题解释：解是否保持物理可行，参数阈值代表什么，近似忽略了哪些机制。

\section{建模流程、量纲与无量纲化}
首先列出可测量的状态及其单位，再用守恒关系或速率关系建立方程，并检查方程是否闭合。例如由 $y'=ry(1-y/K)-H$ 描述持续收获，$y,K$ 的单位是数量，$r$ 为时间的倒数，$H$ 为数量除以时间。取 $u=y/K$、$\tau=rt$，方程化为
\begin{equation}
 \frac{\mathrm du}{\mathrm d\tau}=u(1-u)-h,\qquad h=\frac{H}{rK}.
 \label{eq:ode-harvest-dimensionless}
\end{equation}
三个有量纲参数在无量纲方程中归为一个组合 $h$。这使不同尺度的系统可以比较，但不表示原参数能够分别由单条无量纲轨道识别。

若利用数据估计参数，例如最小化 $\sum_i(y(t_i;\theta)-z_i)^2$，必须说明观测时刻 $t_i$、观测值 $z_i$、未知参数 $\theta$ 的约束和误差假设。目标函数可能有多个极小点；数值拟合结果不等于参数唯一可识别的证明。

\section{种群增长、收获与阈值现象}
设 $H\geq0$。由式\eqref{eq:ode-harvest-dimensionless}，平衡点满足 $u^2-u+h=0$。当 $0<h<1/4$ 时有
\[
 u_\pm=\frac{1\pm\sqrt{1-4h}}2.
\]
导数 $1-2u$ 表明较小平衡点排斥、较大平衡点吸引。$u_-<u_0<u_+$ 的解增加，$u_0>u_+$ 的解减少，均趋于 $u_+$；$0\leq u_0<u_-$ 的解减少并在有限时间到达零。到达零后原常量收获方程会继续产生负数量，故生物解释必须终止或修改收获机制。

当 $h=1/4$ 时，$u'=-(u-1/2)^2$：从上方出发趋于 $1/2$，从下方出发远离它，是单侧吸引的退化平衡。$h>1/4$ 时右端始终小于一个负常数，任意正初值都在有限时间到达零。阈值为 $H_c=rK/4$。两个平衡点在阈值合并并消失，是鞍结分岔（saddle-node bifurcation）的典型现象；这一结论依赖持续、状态无关收获的假设。

\section{机械振动与电路模型}
对 $my''+cy'+ky=F(t)$，取 $\omega_0=\sqrt{k/m}$、$\tau=\omega_0t$、$u=y/Y$，其中 $Y>0$ 为位移尺度，得到
\[
 u_{\tau\tau}+2\zeta u_\tau+u=\frac{F(\tau/\omega_0)}{kY},\qquad
 \zeta=\frac{c}{2\sqrt{mk}}.
\]
阻尼比 $\zeta$ 决定欠阻尼、临界与过阻尼。无外力时能量 $E=m(y')^2/2+ky^2/2$ 满足 $E'=-c(y')^2$，这为解析与数值解提供一致性检查。

串联 RLC 电路以电荷 $q$ 为状态满足
\[
 Lq''+Rq'+q/C=V(t),
\]
其中 $L,C>0$ 为电感与电容，$R\geq0$ 为电阻，电流为 $q'$。它与振子对应 $m\leftrightarrow L$、$c\leftrightarrow R$、$k\leftrightarrow1/C$。电荷和位移虽具有不同单位，数学结构相同。线性元件、理想集中参数等假设仍须分别检验。

\section{捕食模型与传染病模型}
Lotka--Volterra 捕食模型为
\[
 x'=x(a-by),\qquad y'=y(-d+cx),\qquad a,b,c,d>0,
\]
其中 $x$ 为被捕食者数量，$y$ 为捕食者数量。正初值解保持正值，内部平衡点为 $(d/c,a/b)$。直接求导可验证
\[
 H(x,y)=cx-d\log x+by-a\log y
\]
沿解不变。$H$ 在该平衡点严格最小，并在靠近坐标轴或趋于无穷时趋于无穷，因此非平衡正轨道保持在紧能量曲线上。该模型产生围绕内部平衡点的周期运动，而不是吸引的极限环；模型若加入自限增长，结论可能改变。

SIR 模型将总人数 $N$ 分为易感者 $S$、感染者 $I$、移出者 $R$。假设封闭人群、均匀混合、永久移出，采用频率依赖传播：
\begin{equation}
 S'=-\beta\frac{SI}{N},\qquad
 I'=\beta\frac{SI}{N}-\gamma I,\qquad
 R'=\gamma I,
 \label{eq:ode-sir}
\end{equation}
其中 $\beta,\gamma>0$ 单位均为时间的倒数，$S(0)+I(0)+R(0)=N$，初值非负。三个方程相加给出总人数守恒，各坐标面上的向量场保证非负性，因此状态始终在紧单纯形中，解前向全局存在。

对于完全易感背景，基本再生数（basic reproduction number）定义为 $\mathcal R_0=\beta/\gamma$；当前感染人数增加的条件是 $S/N>1/\mathcal R_0$。这说明 $\mathcal R_0>1$ 与给定初始状态下 $I'(0)>0$ 并不完全相同。

若 $S(0)>0$，由 $S'/S=-(\beta/(\gamma N))R'$ 得
\[
 \log\frac{S(t)}{S(0)}=-\frac{\beta}{\gamma N}[R(t)-R(0)].
\]
$R$ 单调有界，故 $\int_0^\infty I(t)\,\mathrm dt<\infty$；$I'$ 有界使 $I$ 一致连续，从而 $I(t)\to0$。否则可抽出高度有正下界、宽度一致的互不相交区间，使积分发散。取极限得到最终规模关系
\[
 \log\frac{S_\infty}{S(0)}
 =-\mathcal R_0\left(1-\frac{S_\infty}{N}-\frac{R(0)}{N}\right).
\]
这些是特定确定性模型的数学结论，参数解释与观测检验仍须依实际数据完成。

\section{综合研究题的组织}
以带收获的 Logistic 模型为例，一个可复核的分析应包含：参数与单位、无量纲化、平衡点及阈值、非负解释的终止条件、数值算法与误差核验。显式解用于核对计算，相线说明全部初值的趋势，稳定性解释扰动，数值实验展示过渡过程；这些方法各自回答不同问题。

对于没有精确解的模型，可用步长减半、不同阶方法比较、守恒量或耗散关系检查数值结果。两种近似一致只是支持证据，不能代替误差界或数学证明。

\section{综合习题与项目}
\begin{enumerate}
 \item\label{ex:ode-harvest} 对 $u'=u(1-u)-h$，分别分析 $h=0,1/8,1/4,1/2$，给出非负初值的长期趋势或到达零的结论。
 \item\label{ex:ode-sir} 在式\eqref{eq:ode-sir}中设 $S(0)/N=0.4$、$\mathcal R_0=2$、$I(0)>0$。判断感染人数初期增加还是减少，并解释。
 \item 证明捕食模型的 $H$ 在内部平衡点严格最小，并说明正能量曲线为何不趋近坐标轴。
 \item 选择收获、振动或 SIR 模型，写出假设、理论分析和误差可复核的数值实验方案。若使用合成数据，应明确标注；若使用观测数据，应注明来源与误差假设。
\end{enumerate}

